\documentclass[border=10pt]{standalone}
\usepackage{fontspec}
\setmainfont{Noto Serif}
\newfontfamily\bengalifont{Noto Serif Bengali}[Script=Bengali]
\usepackage[Latin,Bengali]{ucharclasses}
\setTransitionsFor{Bengali}{\begingroup\bengalifont}{\endgroup}
\usepackage{newunicodechar}
\newunicodechar{।}{{\bengalifont\char"0964}}
\newunicodechar{॥}{{\bengalifont\char"0965}}

\usepackage[table]{xcolor}
\usepackage{booktabs,tabularx,array}
\definecolor{bapilrule}{HTML}{1F2937}
\definecolor{bapilstripe}{HTML}{F3F4F6}
\renewcommand{\arraystretch}{1.28}
\setlength{\tabcolsep}{5pt}

\begin{document}
\begin{minipage}{17cm}
{\large\bfseries সারণি ৪৩।}\ {\large বাংলাদেশের কাঠামোর আলোকে সম্মতি-চেকলিস্ট}\par\vspace{7pt}
\small
\rowcolors{2}{bapilstripe}{white}
\begin{tabularx}{\linewidth}{@{}>{\hsize=0.800\hsize\raggedright\arraybackslash}X>{\hsize=1.224\hsize\raggedright\arraybackslash}X>{\hsize=1.129\hsize\raggedright\arraybackslash}X>{\hsize=0.847\hsize\raggedright\arraybackslash}X@{}}
\arrayrulecolor{bapilrule}\toprule
\rowcolor{white}\textbf{ক্ষেত্র} & \textbf{প্রয়োজনীয়তা (প্রতিবেদিত)} & \textbf{BAPIL নিয়ন্ত্রণ} & \textbf{প্রমাণ} \\
\midrule
সম্মতি ও আইনগত ভিত্তি (PDPA ২০২৬) & ব্যক্তিগত ডেটা সংগ্রহ ও ব্যবহারে সুস্পষ্ট সম্মতি & প্রথমে P0; P3 পাইলট-ডেটার সম্মতির রেকর্ড & সম্মতিপত্র; DPIA \\
ডেটা-বিষয়ীর অধিকার (PDPA ২০২৬) & প্রবেশ, সংশোধন, অপসারণ; স্বয়ংক্রিয় সিদ্ধান্ত সীমিতকরণ & ব্যক্তিসংক্রান্ত কোনো স্বয়ংক্রিয় সিদ্ধান্ত নয়; অধিকার-অনুরোধের প্রক্রিয়া & প্রক্রিয়া; অনুরোধ-লগ \\
সংবেদনশীল ও শিশুদের ডেটা (PDPA ২০২৬) & বাড়তি সুরক্ষা; শিশুদের প্রোফাইলিংয়ে নিষেধাজ্ঞা & সনদ দ্বারা P3 ডেটা বর্জিত & সনদ; ডেটা-তালিকা \\
ডেটা রক্ষণ ও বিনিময় (NDMA ২০২৬) & রক্ষকের দায়িত্ব; অনুমোদিত ডেটা বিনিময় & ডেটা-বিনিময় চুক্তি; উপলব্ধ হলে NRDEX ব্যবহার & স্বাক্ষরিত চুক্তি \\
ডেটা দেশে রাখা (PDPA সংশোধনী, ২০২৬) & নির্দিষ্ট শ্রেণির ডেটা বাংলাদেশে রাখতে হবে (আইনের অধীনে পরিধি ও অবস্থা নিশ্চিত করতে হবে) & P0-এর ঊর্ধ্বের সব ডেটা বাংলাদেশে হোস্টকৃত (NFR-02) & হোস্টিং চুক্তি \\
সোর্স কোড (NDMA ২০২৬) & সরকারি সফটওয়্যারের জাতীয় সোর্স কোড ভান্ডার (২০২৫ সালের অধ্যাদেশের বর্ণনা অনুযায়ী) & উন্মুক্ত সোর্স কোড; উপলব্ধ হলে জাতীয় ভান্ডারে প্রতিচ্ছবি & রিপোজিটরির লিংক \\
নিরাপত্তা-দায়িত্ব (সাইবার নিরাপত্তা আইন ২০২৬) & ব্যবস্থা ও ডেটার সুরক্ষা; ঘটনা ব্যবস্থাপনা & NIST CSF 2.0 নিয়ন্ত্রণ; ঘটনা-পরিকল্পনা & নিরাপত্তা পর্যালোচনা; মহড়া \\
ক্ষতিকর এআই-সৃষ্ট কনটেন্ট (২০২৬ সংশোধনী) & ক্ষতিকর কৃত্রিম কনটেন্ট অপসারণ & কৃত্রিম মিডিয়া তৈরি নয়; P6-এর জন্য গবেষণা-প্রটোকল & ERRC অনুমোদন \\
স্বতঃপ্রণোদিত প্রকাশ (তথ্য অধিকার আইন ২০০৯) & জনগুরুত্বপূর্ণ তথ্য প্রকাশ & উন্মুক্ত প্রকাশনা; স্বচ্ছতা নিবন্ধন & প্রকাশিত নিবন্ধন \\
খসড়া জাতীয় এআই নীতি & ঝুঁকির শ্রেণি; প্রভাব মূল্যায়ন; তত্ত্বাবধান & সামঞ্জস্যপূর্ণ স্তরব্যবস্থা; T2 ও T3-এ প্রভাব মূল্যায়ন & মূল্যায়ন; নিবন্ধন \\
\bottomrule
\end{tabularx}
\par\vspace{5pt}{\footnotesize\itshape প্রয়োজনীয়তাগুলো সংবাদমাধ্যমের প্রতিবেদন ও মাধ্যমিক আইনি ভাষ্যের সারাংশ [84, 85, 86, 87, 88, 113, 114, 115]। কোন বিধান কবে থেকে বলবৎ, তা নিশ্চিত হয়নি। আইনি পর্যালোচনা (অনুচ্ছেদ ১৪.৩) প্রতিটি সারি গেজেটভুক্ত পাঠের সঙ্গে মিলিয়ে দেখবে; ততদিন BAPIL কঠোরতর ব্যাখ্যা প্রয়োগ করবে।}
\par\vspace{3pt}{\scriptsize\color{gray} উৎস: BAPIL কারিগরি ব্লুপ্রিন্ট, সংস্করণ ১.০ (২০২৬)। CC BY 4.0।}
\end{minipage}
\end{document}
